\chapter{Proposal Feedback and Iteration Plan}

\section{Proposal Feedback Summary}

The proposal was defended before the evaluation panel, which returned eleven
substantive comments. Taken together the comments were not a rejection of the
problem or of the architecture; they asked for depth in four specific places:
the project's treatment of language models had to be made explicit rather than
implied, the system had to demonstrate traceability down to the level of code
rather than stopping at artifact relationships, the construction of our
evaluation data had to be defensible rather than asserted, and the dependence
of the design on a human operator had to be reduced and then measured.

The comments and their origin are listed below. Two were raised by more than one
panel member and are recorded once.

\begin{longtable}{p{0.05\textwidth}p{0.52\textwidth}p{0.33\textwidth}}
\toprule
No. & Comment as received & Theme \\
\midrule
\endhead
1  & Research how the involvement of the project manager in the loop can be reduced. & Human-in-the-loop cost \\
2  & Do more research on what kind of language model and what resources will be used. & Model and resource selection \\
3  & Add a traceability aspect to the project. & Traceability \\
4  & Discuss where language models are used, to what capacity or level, and which models are required. & Transparency of model use \\
5  & Dataset validation should also be done. & Evaluation data \\
6  & Code comments carry semantics; take projects with proper code comments. & Retrieval signal and scope \\
7  & Scope should be limited to projects with informative names, comments, directory structure and filenames. & Scope eligibility \\
8  & Robust context management should be factored in. & Context engineering \\
9  & Explore how to use the language model for large code design. & Scalability \\
10 & Include the validation process. & Validation \\
11 & Can functional and non-functional requirements be linked to exact lines of code? & Anchoring granularity \\
\bottomrule
\end{longtable}

Our position on the feedback is that every comment is accepted. Six enter the
committed scope exactly as the panel stated them. Four are accepted with an
explicit stated limit, because the comment as phrased has no natural boundary
and an unbounded commitment cannot be defended at the final evaluation; in each
of those cases the limit is published and measurable rather than left to
judgement. One is deferred to the second phase of the project, contingent on the
committed core being stable.

Accepting this much new work required removing work of comparable size. Two
items left the committed scope to pay for the additions, and the scope changes
are recorded in Section~\ref{sec:ch3-progress} rather than applied silently.
No extension of the project timeline was requested.

\section{Proposal Action Register}

The register below is the authoritative record of what each comment changed. The
evidence column names the artefact that demonstrates the action, so that a
reviewer can verify a claim rather than take it on trust. Status is reported as
of \textbf{[DATE]} and uses three values: \emph{Complete} where the named
evidence exists and has been reviewed by the two members who did not produce it,
\emph{In progress} where work has started and partial evidence exists, and
\emph{Planned} where the design decision is made and recorded but implementation
has not begun.

% ---------------------------------------------------------------------------
% NOTE: Status column is filled with the honest current state. Update the
% [DATE] above and any row whose status changes before submission.
% ---------------------------------------------------------------------------

\begin{longtable}{p{0.03\textwidth}p{0.21\textwidth}p{0.33\textwidth}p{0.10\textwidth}p{0.21\textwidth}}
\toprule
No. & Feedback & Action taken & Status & Evidence \\
\midrule
\endhead
1 & Reduce project-manager involvement
  & Configuration automated and demoted to approval, using state maps and grace
    periods mined from the project's own history. Confirmation narrowed to a
    confidence band and routed to the code owner. Resolution deliberately kept
    human.
  & In progress
  & Use cases and actor split, Ch.~4; intervention metrics in the NFRs \\
2 & Research which model and what resources
  & Selection protocol and candidate matrix written before any model was
    downloaded. No model named on reputation; the protocol fixes the split,
    schema, metrics and thresholds, and the choice is recorded with its version
    and date.
  & In progress
  & Model-selection protocol and candidate matrix, Ch.~5 \\
3 & Add a traceability aspect
  & New \texttt{Requirement} entity and \texttt{SATISFIES} relation carrying a
    commit-pinned \texttt{CodeAnchor}, with re-anchoring by symbol, then content
    hash, then diff hunk. A fourth policy detects requirement--implementation
    divergence.
  & In progress
  & Anchoring diagram and domain model, Ch.~5; policy P4 \\
4 & State where models are used, at what level
  & Usage register added, declaring an autonomy level and a deterministic
    fallback per stage. No model output writes to the graph of truth or fires a
    policy alone; autonomous write is prohibited.
  & Complete
  & Autonomy ladder figure and register table, Ch.~5 \\
5 & Validate the dataset
  & Seven-phase construction and validation protocol defined. Two members label
    each case independently and the third adjudicates; agreement is reported
    against a stated threshold, below which the guide is rewritten and the set
    relabelled.
  & In progress
  & Annotation guide; two annotation files and agreement report, Ch.~8 \\
6 & Code comments carry semantics
  & Each symbol produces a symbol card combining name, signature, docstring,
    comments, path tokens, imports, callers and commit message. The card, not
    the raw line range, is what is embedded and retrieved.
  & Planned
  & Symbol-card design, Ch.~5; comments-stripped ablation, Ch.~8 \\
7 & Limit scope to well-documented projects
  & Accepted with a stated limit. Eligibility score over comment coverage,
    identifier informativeness, directory conventionality, documentation and
    tests; threshold published. A low-documentation cohort is retained and
    reported separately.
  & In progress
  & Eligibility rubric and scoring script; repository catalogue \\
8 & Robust context management
  & Deterministic context packet under a fixed token budget, with per-slot
    allocation, fixed eviction order, syntax-tree-aligned chunking and hybrid
    retrieval. A manifest is stored with every verdict.
  & Planned
  & Context-packet design and verifier-loop diagram, Ch.~5 \\
9 & Use the model for large code design
  & Accepted with a stated limit. Map-then-zoom over a repository map and symbol
    index, full detail only for selected regions, under a fixed pass cap.
    Committed limit: two languages, approximately 500k lines.
  & Planned
  & Scalability benchmark design, Ch.~8 \\
10 & Include the validation process
  & Five-layer plan covering data, component, system, comparative and human
    validation, each with a named gate. Five ablations committed, plus two hard
    gates on uncited verdicts and verdict stability.
  & In progress
  & Validation and acceptance plan, Ch.~8 \\
11 & Link requirements to exact lines of code
  & Accepted with an explicit distinction. Functional requirements anchor to a
    symbol and line range, commit-pinned. Non-functional requirements are not
    line-anchored; they link to an enforcement point and a verifying test, and
    are scored by coverage.
  & In progress
  & Anchoring granularity table, Ch.~4; anchoring diagram, Ch.~5 \\
\bottomrule
\end{longtable}

\section{Iteration Strategy}

Development proceeds in four iterations across two semesters. The strategy is
not iterative merely in the sense of building the product in instalments; each
iteration is organised so that it terminates in something measurable rather than
something merely built, because the central claims of this project are empirical
and an iteration that produces no evidence cannot be assessed.

Three principles govern the sequencing.

\textbf{Risk is retired before it is built upon.} The components whose
feasibility is least certain are scheduled earliest, even where that disturbs a
natural build order. Running an open-weight model on the hardware actually
available to the team, and confirming the event coverage each vendor interface
provides, are both scheduled in the first iteration. Both are external
constraints that cannot be engineered around, and discovering either late would
invalidate work layered on top of it.

\textbf{Deterministic foundations precede probabilistic components.} The event
layer, the version store, the evidence graph and the deterministic rules are
built before the verification layer that consumes them. This ordering follows
from the architecture rather than from preference: the verifier reads evidence
that the deterministic layers produce, and a verifier evaluated against an
unreliable graph measures the graph, not the verifier.

\textbf{Evaluation infrastructure is treated as a deliverable, not as a final
step.} The labelled evaluation set is constructed during the second iteration,
while the components it will eventually judge are still being written. Building
the ground truth before the system that will be measured against it prevents a
subtle failure mode in which cases are unconsciously selected to suit whatever
the implementation happens to handle.

Each iteration closes at a gate. A gate is passed when the named evidence exists
and has been reviewed by the two members who did not produce it, not when the
owning member considers the work finished.

\section{Iteration 1 Plan}

\textbf{Objective.} Establish the foundations on which every later component
depends, and retire the two external unknowns that could still change the
architecture.

\textbf{Selected features.} Requirements and design baseline; read-only
authenticated access to the three source systems; the relationship data model;
deterministic linking from explicit references; local model serving with a
technical spike.

\textbf{Deliverables.} A reproducible repository scaffold with continuous
integration and a container composition; the shared contracts that later work
is written against, namely the normalised event envelope, the node and edge
types and the verdict schema; one event ingested end to end and persisted in the
normalised form; deterministic links resolved on a fixture repository; and a
record of local inference latency, memory and quantisation on the hardware
available to the team.

\textbf{Evidence.} The contract definitions committed as code rather than
described in prose; a passing end-to-end ingestion test; the recorded hardware
measurements; the first set of labelled pilot cases with a measured agreement
figure.

\textbf{Gate.} The team can state, with numbers rather than estimates, what
local inference costs on its own hardware, which vendor events are available and
which are not, and that one event traverses the pipeline and arrives in the
graph in the normalised form.

\section{Iteration 2 Plan}

\textbf{Objective.} Make ingestion survive the delivery behaviour that vendor
interfaces actually exhibit, and construct the evidence on which verification
will later operate.

\textbf{Selected features.} Idempotency, checkpointing and reconciliation;
manual linking and the reviewer-facing dashboard; the code analyser for the
declared languages, producing changed files, symbols and test evidence;
criterion decomposition; the deterministic context-packet builder; and the
labelled evaluation set.

\textbf{Deliverables.} An ingestion path that converges to the same final state
under duplicate and out-of-order delivery and that recovers an event the vendor
never delivered; a populated evidence graph with provenance on every edge; the
analyser producing symbol and test evidence for a labelled pull request; a
deterministically assembled context packet with a recorded truncation manifest;
and the frozen pilot evaluation set with its annotation guide.

\textbf{Evidence.} Controlled event-replay tests demonstrating convergence; the
analyser's extracted output compared against the expected output for a labelled
case; byte-identical packets produced from identical inputs; two independent
annotation files, the agreement statistic computed from them, and a documented
adjudication.

\textbf{Gate.} Reliability is demonstrated rather than asserted, and the
evaluation set is frozen before any component is tuned against it.

\section{Work Breakdown and Responsibilities}

Responsibility is divided by technical stream rather than by architectural
layer, so that each member owns a coherent problem together with the evidence
that demonstrates it works. The test applied to the division was whether
removing any one member would remove a meaningful technical responsibility and
its independent evaluation; a division in which one member owned only an
interface would fail that test.

% ---------------------------------------------------------------------------
% NOTE: this table follows proposal Section 16. The Gantt chart and the
% defence slide assign three of these rows differently. Settle which is
% current and make all three documents agree before submission.
% ---------------------------------------------------------------------------

\begin{longtable}{p{0.16\textwidth}p{0.42\textwidth}p{0.34\textwidth}}
\toprule
Member & Primary ownership & Independent evaluation owned \\
\midrule
\endhead
Haaniah Ismail
  & Source adapters and authenticated read-only access; the normalised event
    envelope; idempotency, checkpointing, reconciliation and stale-source
    handling; the immutable artifact and version store.
  & Duplicate and out-of-order convergence; freshness and recovery after a
    simulated outage; ingestion latency; permission and failure behaviour. \\
Mishal Fatima
  & The code and test analyser for the declared languages; the evidence graph
    and its provenance model; deterministic and manual linking; the
    deterministic drift and policy rules; criterion invalidation and incremental
    dispatch.
  & Link precision and recall; graph correctness; equivalence between
    incremental and full evaluation; rule precision and recall. \\
Abdullah Amin
  & The context-packet builder; local model serving and selection; the bounded
    verifier loop; the parameter-efficient adaptation experiment; the verdict
    and provenance gate; the reviewer evidence report and notification routing.
  & Retrieval, verdict and grounding metrics; the base-versus-adapted
    comparison; local latency and memory; the task-based user evaluation. \\
\bottomrule
\end{longtable}

Four responsibilities are shared, because they cross every stream and cannot be
owned by one member without distorting the result: maintaining one consistent
terminology across the proposal, the report and the implementation;
constructing and adjudicating the labelled evaluation set; end-to-end
integration testing; and the documentation of ethical, privacy and tool-use
disclosures. Within the evaluation set specifically, the member who constructs a
case does not label it, and no member validates the data belonging to their own
subsystem.

\section{Progress Against Plan}
\label{sec:ch3-progress}

% ===========================================================================
% TO COMPLETE BEFORE SUBMISSION.
% Fill from real evidence, not from memory. One row per planned milestone.
% Keep honest statuses: a register in which every row is complete three weeks
% before the evaluation is less credible than one with accurate statuses.
% ===========================================================================

Progress is reported against the plan baselined at the proposal stage, as of
\textbf{[DATE]}. Planned dates are those recorded in the revised schedule;
actual dates are taken from the repository history and the task record rather
than reconstructed afterwards.

\begin{longtable}{p{0.30\textwidth}p{0.14\textwidth}p{0.14\textwidth}p{0.12\textwidth}p{0.22\textwidth}}
\toprule
Milestone & Planned & Actual & Status & Evidence \\
\midrule
\endhead
Proposal revision and panel response & & & & \\
Diagram set (thirteen diagrams) & & & & \\
Repository scaffold, CI, container composition & & & & \\
Shared contracts frozen & & & & \\
Model-selection protocol & & & & \\
Technical spike: local inference measured & & & & \\
Vendor access confirmed & & & & \\
Ingestion: normalised envelope end to end & & & & \\
Evidence graph schema and deterministic linking & & & & \\
Pilot evaluation set constructed & & & & \\
Annotation round and agreement measured & & & & \\
\bottomrule
\end{longtable}

\paragraph{Summary.} [Two or three sentences: how many milestones met on
schedule, how many late, and whether the iteration objective was achieved. State
the overall position plainly rather than presenting the table alone.]

\paragraph{Scope changes since the proposal.} Two items left the committed
scope to pay for the additions accepted in the action register. Narrative
summarisation of change chains was removed: it was already a stretch item, it
produces no evaluable claim, and it is the first element a panel would identify
as decorative. Support for languages beyond the two declared families was
removed, because each additional language multiplies grammar work, convention
modelling and labelling effort without adding a contribution. Additional source
integrations remain outside the committed scope; the adapter layer is defined by
source role rather than by product, so any further source is one adapter written
against an existing contract rather than new architecture, and none produces a
new evaluable claim. [Confirm this paragraph still matches the agreed position
before submission.]

\section{Delays, Deviations and Recovery Plan}

% ===========================================================================
% TO COMPLETE BEFORE SUBMISSION.
% One row per real deviation. Do not invent deviations, and do not omit one
% that is visible in the repository history -- an unexplained gap in the
% commit record is worse than a declared delay.
% ===========================================================================

\begin{longtable}{p{0.22\textwidth}p{0.26\textwidth}p{0.24\textwidth}p{0.22\textwidth}}
\toprule
Deviation & Cause & Impact & Recovery action \\
\midrule
\endhead
 & & & \\
 & & & \\
 & & & \\
\bottomrule
\end{longtable}

\paragraph{Deviations accepted rather than recovered.} [Where a deviation is
not being recovered because the change was deliberate, state that explicitly
here and give the reasoning. A deliberate re-plan is a legitimate outcome; an
undeclared one is not.]

\paragraph{Risk carried into the next iteration.} [Name the risks that remain
open, with the mitigation already prepared for each. The three carried from the
proposal risk register are: vendor interface access or rate limits; criterion to
code retrieval recall falling below the level needed for useful verification;
and local inference latency or memory exceeding the available hardware. Confirm
whether each is still open as of the reporting date.]

\paragraph{Change control.} Every change recorded in this chapter is subject to
the change-control rule adopted at the proposal stage: each change is recorded
with its date, the evidence motivating it, the approver, the effect on the
schedule and the impact on traceability. Optional features do not enter the
committed scope merely because routine implementation finishes ahead of
schedule.
